\section{Introducción}
\label{sec:introduccion}

\pendiente{Objetivo del trabajo y organización del informe.}

\subsection{Electrocardiograma y derivaciones}
\label{sec:derivaciones}
% Responde: ítems 3b y 3c
\pendiente{Ítem 3b: por qué en los estudios clínicos se colocan electrodos en miembros superiores e inferiores.}

\pendiente{Ítem 3c: otras derivaciones del ECG y qué información espacial de la actividad eléctrica cardíaca aportan al profesional.}

\subsection{Formatos de almacenamiento de ECG}
\label{sec:formatos}
% Responde: ítem 1 (segunda parte)
\pendiente{Ítem 1: otros formatos en los que se registran electrocardiogramas.}

\subsection{Variabilidad de la frecuencia cardíaca}
\label{sec:hrv}
% Base teórica para los ítems 6, 7 y 8
\pendiente{Marco teórico de la VFC necesario para interpretar el tacograma y el periodograma; bandas de análisis con su fuente.}
